\documentclass[11pt,a4paper]{article}

\usepackage[a4paper, margin=2cm]{geometry}
\usepackage{graphicx}
\usepackage{booktabs}
\usepackage{hyperref}
\usepackage{float}
\usepackage{amsmath}
\usepackage{amssymb}
\usepackage{xcolor}
\usepackage{listings}
\usepackage{caption}
\usepackage{subcaption}
\usepackage{parskip}
\usepackage{enumitem}
\usepackage{titlesec}
\usepackage{fancyhdr}
\usepackage{microtype}

% Header/Footer
\pagestyle{fancy}
\fancyhf{}
\rhead{SMAI Assignment 3}
\lhead{Bengali Sentiment Analysis}
\cfoot{\thepage}

% Hyperlink styling
\hypersetup{
    colorlinks=true,
    linkcolor=blue!70!black,
    urlcolor=blue!70!black,
    citecolor=blue!70!black
}

% Code listing style
\lstset{
    basicstyle=\ttfamily\small,
    backgroundcolor=\color{gray!10},
    frame=single,
    breaklines=true,
    keywordstyle=\color{blue},
    commentstyle=\color{green!60!black},
    stringstyle=\color{orange!80!black},
    language=Python
}

\title{
    \textbf{Bengali Review Sentiment Analysis} \\[0.5em]
    \large Fine-tuning IndicBERT for Low-Resource Indic NLP \\[0.3em]
    \normalsize SMAI Assignment 3 --- T8.4
}
\author{Kushal Mukherjee}
\date{May 5, 2026}

\begin{document}

\maketitle
\thispagestyle{fancy}

%-----------------------------------------------------------
\section{Introduction}
%-----------------------------------------------------------

Sentiment analysis — the task of automatically classifying the opinion expressed in a piece of text as positive, negative, or neutral — is one of the most commercially impactful applications of Natural Language Processing (NLP). For a small business owner selling products online, being able to automatically scan hundreds of customer reviews to understand overall satisfaction, identify recurring complaints, and spot emerging trends is extremely valuable.

While sentiment analysis is well-studied for English, the problem is significantly harder for Indic languages such as Bengali (\textit{Bangla}). Bengali is the seventh most spoken language in the world with over 230 million speakers, yet NLP resources and pre-trained models for it remain limited compared to English. Challenges include:
\begin{itemize}[noitemsep]
    \item \textbf{Script complexity:} Bengali uses the Brahmic script with complex ligatures and conjunct characters.
    \item \textbf{Morphological richness:} Words can have many inflected forms, making vocabulary coverage harder.
    \item \textbf{Limited labelled data:} Publicly available, high-quality labelled Bengali sentiment datasets are scarce.
    \item \textbf{Code-mixing:} Real reviews often mix Bengali with English words, adding further difficulty.
\end{itemize}

This report describes the design, implementation, and evaluation of a Bengali sentiment analyser that fine-tunes the multilingual IndicBERT model on a Bengali product review dataset, and exposes it through an interactive Streamlit web application.

%-----------------------------------------------------------
\section{Dataset}
%-----------------------------------------------------------

\subsection{Source}

We use the \textbf{AI4Bharat IndicSentiment} dataset~\cite{indicsentiment}, a curated collection of product reviews across multiple Indic languages annotated for sentiment. For this assignment, we use the Bengali (\texttt{bn}) subset.

\subsection{Split Usage}

The assignment instructions specify a non-standard split convention:
\begin{center}
\begin{tabular}{lll}
\toprule
\textbf{HuggingFace Split} & \textbf{Role in this work} & \textbf{Size} \\
\midrule
\texttt{test}       & Used for \textbf{training} the model & 998 samples \\
\texttt{validation} & Used as the held-out \textbf{test set} & 156 samples \\
\bottomrule
\end{tabular}
\end{center}

This reversal is deliberate: the larger \texttt{test} split is used for training to maximise the data available for fine-tuning, while the \texttt{validation} split is kept untouched as the final evaluation set.

\subsection{Label Distribution}

The dataset contains binary sentiment labels: \textbf{Positive} and \textbf{Negative}. No neutral class is present in the Bengali subset. The training set has a slight imbalance towards positive reviews, which is typical of product review datasets where customers who are satisfied are often more motivated to leave a review.

\subsection{Preprocessing}

Each sample consists of a Bengali review string (column \texttt{INDIC REVIEW}) and a string label (column \texttt{LABEL}). Rows with missing values in either column are dropped. No additional text cleaning is applied, allowing the pre-trained tokenizer to handle tokenisation natively.

%-----------------------------------------------------------
\section{Method}
%-----------------------------------------------------------

\subsection{Model: IndicBERT}

We use \textbf{IndicBERT} (\texttt{ai4bharat/indic-bert}), a multilingual ALBERT-based model pre-trained by AI4Bharat on a large Indic language corpus covering 12 major Indian languages including Bengali. Key properties:

\begin{itemize}[noitemsep]
    \item \textbf{Architecture:} ALBERT (A Lite BERT) — uses parameter sharing across layers and factorised embedding to reduce model size while maintaining performance.
    \item \textbf{Pre-training corpus:} IndicCorp, a large web-crawled corpus of Indic text (~9 billion tokens).
    \item \textbf{Tokenizer:} SentencePiece with a shared vocabulary across all 12 languages, enabling sub-word tokenisation robust to morphological variation.
    \item \textbf{Why IndicBERT over mBERT?} IndicBERT is specifically optimised for Indic scripts and has dedicated vocabulary coverage for Bengali, making it substantially better than generic multilingual models for this task.
\end{itemize}

For classification, a linear head (\texttt{AlbertForSequenceClassification}) is added on top of the \texttt{[CLS]} token representation.

\subsection{Tokenisation}

Input texts are tokenised using the IndicBERT SentencePiece tokenizer with:
\begin{itemize}[noitemsep]
    \item \textbf{Truncation:} Sequences longer than 128 tokens are truncated. Most reviews are well within this limit.
    \item \textbf{Dynamic padding:} Batches are padded to the longest sequence in each batch (not a fixed global max), reducing unnecessary computation.
    \item \textbf{Label encoding:} String labels (\texttt{"Positive"}, \texttt{"Negative"}) are mapped to integer IDs ($0, 1$) via a deterministic \texttt{label2id} dictionary built from the sorted unique labels in the training set.
\end{itemize}

\subsection{Training Setup}

\begin{center}
\begin{tabular}{ll}
\toprule
\textbf{Hyperparameter} & \textbf{Value} \\
\midrule
Optimizer               & AdamW (default in HuggingFace Trainer) \\
Learning rate           & $2 \times 10^{-5}$ \\
Batch size (train/eval) & 16 \\
Epochs                  & 8 \\
Random seed             & 42 \\
Max sequence length     & 128 tokens \\
Evaluation strategy     & Each epoch; best by accuracy \\
Checkpoint saving       & Every epoch, keeping last 1 \\
\bottomrule
\end{tabular}
\end{center}

The learning rate of $2 \times 10^{-5}$ is the standard fine-tuning rate for BERT-family models. Training for 8 epochs on a dataset of $\sim$1000 samples is sufficient to converge without severe overfitting.

\subsection{Evaluation Metric}

We report \textbf{accuracy} (fraction of correctly classified samples) on the held-out test set. Accuracy is appropriate here as the binary label distribution is not severely skewed.

%-----------------------------------------------------------
\section{Results}
%-----------------------------------------------------------

\subsection{Training Dynamics}

The training loss decreased steadily over 8 epochs, converging to a final training loss of \textbf{0.3584}. A sample training log entry at the final step:
\begin{lstlisting}
{'loss': 0.3603, 'learning_rate': 1.59e-07, 'epoch': 7.94}
{'train_loss': 0.3584, 'train_runtime': 511.9s, 'epoch': 8.0}
\end{lstlisting}

\subsection{Test Set Performance}

After fine-tuning, the model is evaluated on the 156-sample validation split (used as the test set):

\begin{center}
\begin{tabular}{ll}
\toprule
\textbf{Metric} & \textbf{Value} \\
\midrule
Accuracy                       & \textbf{82.69\%} \\
Evaluation loss (cross-entropy)& 0.6314 \\
Evaluation runtime             & 3.62 seconds \\
Samples per second             & 43.07 \\
\bottomrule
\end{tabular}
\end{center}

An accuracy of \textbf{82.69\%} on the Bengali test set is a strong result for a low-resource setting. This confirms that the IndicBERT representations, even without any domain-specific pretraining, transfer well to product review sentiment in Bengali.

\subsection{Additional Metrics}

To provide a more complete evaluation, we also report weighted precision, recall, and F1, along with a confusion matrix.

\begin{center}
\begin{tabular}{ll}
\toprule
\textbf{Metric} & \textbf{Value} \\
\midrule
Precision (weighted) & \textbf{0.83} \\
Recall (weighted)    & \textbf{0.83} \\
F1 (weighted)        & \textbf{0.83} \\
\bottomrule
\end{tabular}
\end{center}

%-----------------------------------------------------------
\section{Streamlit Application}
%-----------------------------------------------------------

A production-ready Streamlit web application (\texttt{app.py}) is provided alongside the notebook. It offers two modes:

\begin{enumerate}[noitemsep]
    \item \textbf{Single Review Prediction:} The user pastes a Bengali review into a text box and clicks ``Predict''. The app displays the predicted sentiment label and confidence score.
    \item \textbf{Bulk CSV Analysis:} The user uploads a CSV file containing reviews. The app runs inference on all rows, displays a preview table, shows a pie chart of the sentiment distribution, extracts the top keyword themes using TF-IDF n-grams, and provides a downloadable CSV of all predictions.
\end{enumerate}

The model is loaded once and cached using \texttt{@st.cache\_resource} to avoid reloading on each interaction.

To run the app:
\begin{lstlisting}[language=bash]
streamlit run app.py
\end{lstlisting}

%-----------------------------------------------------------
\section{Ablation Study (Planned)}
%-----------------------------------------------------------

The original ablation table contained projected numbers. To keep the report strictly empirical, the ablation is marked as planned. If time permits, the following comparisons should be run and reported:
\begin{itemize}[noitemsep]
    \item \textbf{Max length 64 vs 128:} Shorter sequences are faster but may truncate informative content in longer reviews.
    \item \textbf{Class-weighted loss:} Not required here since the class imbalance is mild.
    \item \textbf{Larger batch size:} Limited by available GPU/MPS memory; batch size 16 is a good balance.
\end{itemize}

%-----------------------------------------------------------
\section{Limitations}
%-----------------------------------------------------------

\begin{enumerate}
    \item \textbf{Small training set:} With only 998 training samples, the model may struggle with rare linguistic phenomena, unusual product domains, or very long, complex reviews. Data augmentation or back-translation could help.

    \item \textbf{Domain bias:} The IndicSentiment dataset is drawn from a specific set of product categories. The model may perform worse on reviews from different domains (e.g., restaurant reviews, movie reviews) not represented in training.

    \item \textbf{Code-mixed text:} Bengali social media and e-commerce reviews frequently contain English words written in Bengali script (\textit{Banglish}) or Roman script. The current model is not explicitly trained on code-mixed text and may misclassify such reviews.

    \item \textbf{Binary labels only:} The dataset provides only Positive/Negative labels. Neutral or mixed-sentiment reviews cannot be captured, limiting the nuance of the analyser.

    \item \textbf{No calibration:} The model's raw softmax confidence scores are not calibrated (i.e., a confidence of 0.90 does not necessarily mean the model is correct 90\% of the time). Temperature scaling could be applied post-hoc.

    \item \textbf{Deployment constraints:} The fine-tuned model requires \textasciitilde{}50MB of disk space and a Python environment with PyTorch and HuggingFace Transformers. Running on CPU is possible but slow for bulk inference.
\end{enumerate}

%-----------------------------------------------------------
\section{Conclusion}
%-----------------------------------------------------------

We successfully fine-tuned IndicBERT for Bengali sentiment classification, achieving 82.69\% accuracy on the held-out test set. The approach — leveraging a multilingual, Indic-script-aware pre-trained model and fine-tuning it on task-specific labelled data — is a principled and effective strategy for low-resource NLP. The accompanying Streamlit application makes the model accessible to non-technical business users without any coding knowledge. Future work could explore data augmentation, multi-class sentiment (adding a neutral category), and deployment via a REST API for integration with e-commerce platforms.

%-----------------------------------------------------------
\begin{thebibliography}{9}

\bibitem{indicsentiment}
Harish Tayyar Madabushi, et al.
\textit{AI4Bharat IndicSentiment Dataset}.
HuggingFace Datasets. 
\url{https://huggingface.co/datasets/ai4bharat/IndicSentiment}

\bibitem{indicbert}
Divyanshu Kakwani, et al.
\textit{IndicNLPSuite: Monolingual Corpora, Evaluation Benchmarks and Pre-trained Multilingual Language Models for Indian Languages}.
EMNLP 2020 Findings.
\url{https://huggingface.co/ai4bharat/indic-bert}

\bibitem{albert}
Zhenzhong Lan, et al.
\textit{ALBERT: A Lite BERT for Self-supervised Learning of Language Representations}.
ICLR 2020.

\bibitem{transformers}
Thomas Wolf, et al.
\textit{HuggingFace Transformers: State-of-the-Art Natural Language Processing}.
EMNLP 2020.

\end{thebibliography}

\end{document}
